\documentclass[pdflatex,sn-mathphys-num]{sn-jnl}

\usepackage{graphicx}
\usepackage{listings}
\usepackage{xcolor}
\usepackage{booktabs}
\usepackage{float}
\usepackage{caption}
\usepackage{subcaption}
\usepackage{amsmath}
\usepackage{tcolorbox}
\usepackage{lipsum}
\usepackage{multirow}
\usepackage{mathrsfs}
\usepackage{amssymb}
\usepackage{algorithm}
\usepackage{xcolor}
\usepackage{xcolor}
\usepackage{listings}
\usepackage{enumitem} 
% Custom color scheme
\definecolor{delim}{RGB}{20,105,176}
\definecolor{numb}{RGB}{106, 115, 125}
\definecolor{string}{RGB}{163, 21, 21}
\definecolor{comment}{RGB}{0, 128, 0}
\definecolor{keyword}{RGB}{0, 0, 255}
\begin{document}

\title{New Generation Data Models and DBMSs Course - Assignment 3 - NoSQL Queries}

\author*[1]{\fnm{Carlo} \sur{Falchi (Matr. 67340A)}}\email{carlo.falchi@studenti.unimi.it}
\author*[1]{\fnm{Francesco} \sur{Palladini (Matr. 85652A)}}\email{francesco.palladini@studenti.unimi.it}
\affil[1]{\orgdiv{Department of Computer Science}, \orgname{University of Milan}, \country{Italy}}
\lstdefinelanguage{json}{
    basicstyle=\small\ttfamily,
    numbers=left,
    numberstyle=\tiny\color{gray},
    stepnumber=1,
    numbersep=8pt,
    showstringspaces=false,
    breaklines=true,
    frame=lines,
    backgroundcolor=\color{white},
    commentstyle=\color{comment}\itshape,
    morecomment=[s]{/*}{*/},
    morecomment=[l]{//},
    stringstyle=\color{string},
    morestring=[b]"
}
\maketitle

\section{Assignment}
\begin{enumerate}
\item For each university, identify the name, kind and city of the university that are located in California or Massachusetts along with the name, cfu, teacher surname and number of participants of the courses that it offers.
\item For each university, identify the teachers that teach the top three courses in terms of number of participants.
\item Identify the students that passed all their exams for courses with a written text with an average mark above 27.
\item For each course, report the list of students that participated to the course. Report also the number of students that passed the course with a mark between 18 and 23, between 24 and 27, and above 27. Report also the number of students that did not pass the course.
\item Identify the pair of students that are friend of a friend (i.e. a student should be associated with the friends of his friends).
\end{enumerate}
\subsection{First Query}
For each university, identify the name, kind and city of the university that are located in California or Massachusetts along with the name, cfu, teacher surname and number of participants of the courses that it offers.
\begin{enumerate}[label=(\alph*)]
\item First, we filter out only the universities in California or Massachusetts
\item The document is then joined according to the "\_id" to associate all the courses of each university
\item The field courses is unwind to get the single element of the array
\item A new field is added to calculate the number of participants for that course
\item The group by operation is needed to group all the university's courses for each university \_id
\end{enumerate}
\begin{lstlisting}[language=json]
db.university.aggregate([
  {
    $match: {
      $or: [
        { state: "California" },
        { state: "Massachusetts" }
      ]
    }
  },
  {
    $lookup: {
      from: "course",
      localField: "_id",
      foreignField: "university",
      as: "courses"
    }
  },
  {
    $unwind: { path: "$courses" }
  },
  {
    $addFields: {
      numPartecipants: { $size: "$courses.participants" }
    }
  },
  {
    $group: {
      _id: "$_id",
      name: { $first: "$name" },
      kind: { $first: "$kind" },
      city: { $first: "$city" },
      courses: { $push: "$courses" }
    }
  }
])
\end{lstlisting}

\subsection{Second Query}
For each university, identify the teachers that teach the top three courses in terms of number of participants.
\begin{enumerate}[label=(\alph*)]
\item First, we add a field that counts the number of participants for every course
\item Then we group by University and Teacher Name and Surname, counting for each one the course with the max number of participants
\item Sort the documents list in descending order
\item We group again by university to limit to the first 3 teacher for each university
\item The document is joined according to the \_id field which is the university id, to associate the university info
\item We got an array with one element so we unwind it
\item Finally, we show only the requested fields: University Name and the related Teachers
\end{enumerate}
\begin{lstlisting}[language=json]
db.course.aggregate([
  { "$addFields": { "numParticipants": { "$size": "$participants" } } },
  { "$group": { "_id": { "univ": "$university", "name": "$teacher.name", "surname": "$teacher.surname" }, "maxParticipants": { "$max": "$numParticipants" } } },
  { "$sort": { "maxParticipants": -1 } },
  { "$group": { "_id": "$_id.univ", "teachers": { "$firstN": { "input": { "name": "$_id.name", "surname": "$_id.surname" }, "n": 3 } } } },
  { "$lookup": { "from": "university", "localField": "_id", "foreignField": "_id", "as": "univ" } },
  { "$unwind": { "path": "$univ" } },
  { "$project": { "name": "$univ.name", "teachers": "$teachers" } }
])
\end{lstlisting}
\subsection{Third Query}
Identify the students that passed all their exams for courses with a written text with an average mark above 27.
\begin{enumerate}[label=(\alph*)]
\item First, we do an unwind of the exams to get one for each document
\item The document is joined with the related course for each exam passed by the student
\item We use a match stage to filter out only the written exams
\item We group by the student to calculate the average of the resulting exams
\item Filter out only the students that have an average above 27 
\item We remove the unrequested fields by project stage
\end{enumerate}
\begin{lstlisting}[language=json]
db.student.aggregate([
  { $unwind: "$exams" },
  {
    $lookup: {
      from: "course",
      localField: "exams.course",
      foreignField: "_id",
      as: "course",
      pipeline: [{ $project: { exam_type: 1 } }]
    }
  },
  { $match: { "course.exam_type": "Written" } },
  {
    $group: {
      _id: "$_id",
      average: { $avg: "$exams.mark" },
      name: { $first: "$name" },
      surname: { $first: "$surname" }
    }
  },
  { $match: { average: { $gt: 27 } } },
  { $project: { _id: 0, average: 0 } }
])
\end{lstlisting}
\subsection{Fourth Query}
For each course, report the list of students that participated to the course. Report also the number of students that passed the course with a mark between 18 and 23, between 24 and 27, and above 27. Report also the number of students that did not pass the course.
\begin{enumerate}[label=(\alph*)]
\item First, we do an unwind of the participants to get an element for the document
\item Then, we join according to the participants field in order to associate the related student info
\item Given that we got an array of students, we do unwind to extract the only object element in the array
\item To keep track only of exam related to the course document within the exams student array, we extract it as new fields using the \$filter expression
\item For each document, we calculate those that respect the requested range of marks.
\item Finally, we group by course to give a sum of the requested numbers.
\end{enumerate}
\begin{lstlisting}[language=json]
db.course.aggregate([
  { $unwind: "$participants" },
  {
    $lookup: {
      from: "student",
      localField: "participants",
      foreignField: "_id",
      as: "student"
    }
  },
  { $unwind: "$student" },
  {
    $addFields: {
      currentCourseExam: {
        $first: {
          $filter: {
            input: "$student.exams",
            as: "exam",
            cond: { $eq: ["$$exam.course", "$_id"] }
          }
        }
      }
    }
  },
  {
    $addFields: {
      didntPassExam: {
        $cond: [{ $not: ["$currentCourseExam"] }, 1, 0]
      },
      passedwith18to23: {
        $cond: [
          {
            $and: [
              { $gte: ["$currentCourseExam.mark", 18] },
              { $lte: ["$currentCourseExam.mark", 23] }
            ]
          },
          1,
          0
        ]
      },
      passedwith24to27: {
        $cond: [
          {
            $and: [
              { $gte: ["$currentCourseExam.mark", 24] },
              { $lte: ["$currentCourseExam.mark", 27] }
            ]
          },
          1,
          0
        ]
      },
      passedabove27: {
        $cond: [{ $gt: ["$currentCourseExam.mark", 27] }, 1, 0]
      }
    }
  },
  {
    $group: {
      _id: "$_id",
      name: { $first: "$name" },
      participants: { $push: "$student" },
      total_passedwith18to23: { $sum: "$passedwith18to23" },
      total_passedwith24to27: { $sum: "$passedwith24to27" },
      total_passedabove27: { $sum: "$passedabove27" },
      total_didntPassExam: { $sum: "$didntPassExam" }
    }
  }
])
\end{lstlisting} 
\subsection{Fifth Query}
Identify the pair of students that are friend of a friend (i.e. a student should be associated with the friends of his friends).
\begin{enumerate}[label=(\alph*)]
\item First, we unwind the list of friends
\item Then, we associate the friend info with the student collection
\item Given that friend\_info is an array, we unwind it to get the only element of the array, and then we unwind the friends of friends
\item We then group by id and the friends of friends to avoid duplicates
\item We join the student collection related to the friend of friend and unwind it
\item Finally, we project the requested fields
\end{enumerate}
\begin{lstlisting}[language=json]
db.student.aggregate(
[
    { $unwind: "$friends" },
    { 
        $lookup: { from: "student", localField: "friends", foreignField: "_id", as: "friend_info" } 
    },
    { $unwind: "$friend_info" },
    { $unwind: "$friend_info.friends" },
    { 
        $group: { 
            _id: { _id: "$_id", friendOfFriendId: "$friend_info.friends" }, 
            name: { $first: "$name" }, 
            surname: { $first: "$surname" } 
        } 
    },
    { 
        $lookup: { from: "student", localField: "_id.friendOfFriendId", foreignField: "_id", as: "friendoffriend" } 
    },
    { $unwind: "$friendoffriend" },
    { $project: { _id: 0, name: 1, surname: 1, friendofFriendName: "$friendoffriend.name", friendofFriendSurname: "$friendoffriend.surname" } }
])    
\end{lstlisting} 

\end{document}